\section{Reproducibility, the detail}
\label{app:repro}

\paragraph{What a fresh clone holds.} The repository contains the study code, the offline tests, the
protocol with its dated additions, the declarations under \texttt{study/} (the 18 comparisons, the run
specification with model revisions, the 20 validation repositories, both SQL queries), and the evidence
package under \texttt{reports/registered-2026-09-13/}: the frozen manifest and aggregate daily signals, the
full eligible frame and its frozen order, the skip log, the query exports and provenance, the hourly
validation comparison and figures, the signed readiness review, the run metadata with package versions and
hardware, the integrity log, the confirmatory inference tables, the declared descriptive tables and the
post-hoc accounting. It does not contain raw endpoint responses, model weights, the version 1 reference copy,
or the private research notes. The BigQuery frame and validation queries billed 12,510,623,170,560 bytes
(about \$71.11 at the list price before credits), approved in writing before they ran.

\paragraph{Commands.} Place the release asset at
\texttt{reports/registered-2026-09-13/run/results.parquet}, then from the repository root, with the
project's Python environment:
{\footnotesize
\begin{verbatim}
.\.venv\Scripts\python.exe -m scripts.review_integrity --run reports/registered-2026-09-13/run
.\.venv\Scripts\python.exe -m scripts.registered_inference --run reports/registered-2026-09-13/run `
    --out data/cache/inference-reproduction
.\.venv\Scripts\python.exe reports/registered-2026-09-13/summarize_descriptives.py
\end{verbatim}}
The first verifies hashes, the grid, contexts, targets, labels, seeds and recomputes every score; the second
regenerates the 18-row inference table from the saved forecasts without refitting; the third regenerates the
descriptive CSVs behind the tables of Appendix~\ref{app:tables}. The post-hoc accounting is regenerated by
\texttt{reports/registered-2026-09-13/posthoc/accounting.py}, and the figures of this report by
\texttt{paper/figures/make\_figures.py}, which reads only the saved CSVs. The run itself used an NVIDIA
GeForce RTX 4070 SUPER (driver 610.47, CUDA 12.8, PyTorch 2.11.0), \texttt{chronos-forecasting} 2.3.2,
\texttt{timesfm} 3.0.2 and \texttt{statsforecast} 2.1.1 \citep{timesfm-pypi, statsforecast-2.1.1}.

\paragraph{Licences.} The code is MIT. The aggregate dataset and this write-up are CC BY 4.0. Chronos-2 and
TimesFM 2.5 weights are Apache-2.0; TimesFM 3 weights carry the TimesFM Non-Commercial License v1.0; the
Chronos and TimesFM packages are Apache-2.0.
